\section{Computed party-permutation extensions}
\label{sec:party-extensions}

This appendix records finite examples of the projective extension in
Corollary~\ref{cor:discrete-lu-symmetry}.  They are not used in the
uniform rigidity or transversal-group theorems.

\begin{corollary}[Computed party-extension splittings]
\label{cor:computed-party-splitting}
The realized projective party-permutation extension splits for both
non-GRS \(q=11\) pencil classes, the \(q=11\) GRS evaluation-set
representatives
\[
\{0,1,2,3,4,5\},\ \{0,1,2,3,4,6\},\
\{0,1,2,3,5,6\},\ \{0,1,2,3,5,9\},
\]
the enhanced-symmetry \(t=-1\) rows over \(\F_{17}\) and \(\F_{31}\),
and the \(H_3\) representatives
\((q,\tau)=(11,4),(19,5),(29,6),(31,13)\).
For every listed non-GRS or \(H_3\) row, party parity acts on the fixed
split torus \(T\) by inversion.  Thus, after choosing an anchor leg and
allowing party motion to transport the corresponding encoder view, a
realized odd party permutation enlarges the possible local linear
carrier exactly from \(T\) to
\[
 N(T)=\langle T,J\rangle,\qquad
 J=\begin{pmatrix}0&-1\\1&0\end{pmatrix},\qquad
 J^2=-I,\qquad J\diag(a,a^{-1})J^{-1}=\diag(a^{-1},a).
\]
Hence \(N(T)/T\cong C_2\).  The linear extension does not split in odd
characteristic, although its projectivization does.  If the party
permutation fixes the input leg of a chosen encoder, this carrier acts
as an actual logical enlargement; otherwise it compares two encoder
views of the same six-leg tensor.
\end{corollary}

\begin{proof}
For each row, an exact finite computation reconstructs the fifteen
four-support shortened planes, exhausts all \(720\) party permutations,
and checks the propagated local symplectic blocks against the full CSS
Lagrangian.  After quotienting by the fixed-party group identified in
Theorem~\ref{thm:logical-phase}, it is enough to test anchor block \(I\)
on the GRS locus and anchor blocks \(I,J\) off it; this discards no
candidate.  In the order in which the six \(q=11\) rows are listed above,
their party images are
\[
 S_5,\quad S_4,\quad V_4,\quad C_5,\quad S_3,\quad D_{12},
\]
where \(D_{12}\) has order \(12\).  The \(q=17\) and \(q=31\)
enhanced rows have party images \(S_4\) and \(S_5\), respectively, and
each listed \(H_3\) image is \(S_5\).

The certificate records the complete normalized factor set for each
section.  It also closes a complement of order \(\lvert\Pi\rvert\) whose
projection to \(\Pi\) is bijective, and verifies that the associated
normalized cochain changes the factor set to the identity.  Hence every
listed extension splits.  Conjugating the fixed torus by the stored lifts
gives the sign action; any odd lift therefore induces the nontrivial
class of \(N(T)/T\).  The complement is an involution in the projective
party-extension group; its determinant-one anchor block is \(J\), whose
square is \(-I\).
\end{proof}

At \(q=11\), the two non-GRS pencil classes have party images \(S_5\) and
\(S_4\), their extensions split, and odd party motion inverts \(T\).
After choosing an anchor leg, the possible one-qudit affine carriers
across transported encoder views therefore form
\(\F_q^2\rtimes N(T)\).  This is the logical group of a fixed encoder
only for the subgroup fixing its input leg; general party motion compares
different encoder views.  The GRS fixed-party logical group and this
anchored non-GRS carrier have orders \(159720\) and \(2420\), respectively.

Arithmetic symmetry jumps do not change the logical phase.  The
characteristic-\(17\) tetrahedral specialization lies on the GRS boundary
and retains the full \(\mathrm{SL}_2(17)\) fixed-party image; its party
image is \(S_4\), and the computed extension splits.  The
characteristic-\(31\) \(A_5\) specialization is non-GRS, so its
fixed-party image is the split torus; its \(S_5\) party image also splits,
and odd motion enlarges the transported anchor-block carrier to \(N(T)\).
